% Adaptación de FORMATO.tex para ACD3; plantilla original preservada.
\documentclass[12pt,a4paper]{report}

% ------------------------------------------
%   PAQUETES BÁSICOS
% ------------------------------------------
\usepackage[spanish]{babel}
\usepackage[utf8]{inputenc}
\usepackage[T1]{fontenc}
\usepackage{graphicx}
\usepackage{amsmath, amssymb}
\usepackage{booktabs}
\usepackage{geometry}
  \geometry{left=3cm, right=2.5cm, top=3cm, bottom=2.5cm}
\usepackage{setspace}
  \onehalfspacing
\usepackage{titlesec}
  \titleformat{\chapter}[block]{\bfseries\large}{\thechapter.}{1em}{}
  \titleformat{\section}[block]{\bfseries}{\thesection.}{1em}{}
\usepackage{tocloft}
  \setlength{\cftbeforechapskip}{6pt}
  \renewcommand{\cftchapfont}{\bfseries}
\usepackage{longtable}
\usepackage{array}
\usepackage{verbatim}

\usepackage{cite}
\usepackage{hyperref}
% Soporte mínimo para tablas, rutas y bloques técnicos.
\usepackage{calc}
\usepackage{xurl}
\usepackage{fvextra}
\providecommand{\tightlist}{\setlength{\itemsep}{0pt}\setlength{\parskip}{0pt}}
\setlength{\emergencystretch}{3em}

\addto\captionsspanish{
    \renewcommand{\contentsname}{Índice General}
    \renewcommand{\tablename}{Tabla}
}

% ------------------------------------------
%   FONDO DE PÁGINA (SOLO CARÁTULA)
% ------------------------------------------
\usepackage{eso-pic}
\newcommand\BackgroundPic{%
  \put(0,0){%
    \parbox[b][\paperheight]{\paperwidth}{%
      \vfill
      \centering
      \IfFileExists{fondo.jpeg}{\includegraphics[width=\paperwidth,height=\paperheight]{fondo.jpeg}}{\parbox{0.8\paperwidth}{\centering\small [PENDIENTE: fondo.jpeg original]}}%
      \vfill
}}}

% ------------------------------------------
%   ENCABEZADO CON LOGO (OTRAS PÁGINAS)
% ------------------------------------------
\usepackage{fancyhdr}
\pagestyle{fancy}
\fancyhf{}
\setlength{\headheight}{40pt}
\fancyhead[L]{\IfFileExists{logo_unl.jpeg}{\includegraphics[height=1.2cm]{logo_unl.jpeg}}{\fbox{\scriptsize PENDIENTE: logo\_unl.jpeg}}}
\fancyfoot[C]{\thepage}
\renewcommand{\headrulewidth}{0pt}

\fancypagestyle{plain}{
    \fancyhf{}
    \fancyhead[L]{\IfFileExists{logo_unl.jpeg}{\includegraphics[height=1.2cm]{logo_unl.jpeg}}{\fbox{\scriptsize PENDIENTE: logo\_unl.jpeg}}}
    \fancyfoot[C]{\thepage}
    \renewcommand{\headrulewidth}{0pt}
}

% ------------------------------------------
%   METADATOS DEL DOCUMENTO
% ------------------------------------------
\newcommand{\TituloTrabajo}{ACD3 -- Clínica de arquitectura y contrato de API}
\newcommand{\Autor}{[PENDIENTE: integrantes del equipo]}
\newcommand{\Director}{[PENDIENTE: docente]}
\newcommand{\Carrera}{Carrera de Computación}
\newcommand{\Facultad}{Facultad de la Energía, las Industrias y los Recursos Naturales No Renovables}
\newcommand{\Universidad}{Universidad Nacional de Loja}
\newcommand{\Ciudad}{Loja, Ecuador}
\newcommand{\Fecha}{Semana 3 -- [PENDIENTE: fecha real]}

\begin{document}

\AddToShipoutPictureBG*{\BackgroundPic}

% ------------------------------------------
%   PORTADA
% ------------------------------------------
\begin{titlepage}
  \centering
  \vspace*{2cm}
  {\bfseries\Large \Universidad}\\[0.5cm]
  {\bfseries\large \Facultad}\\[0.5cm]
  {\bfseries\large \Carrera}\\[2cm]
  {\bfseries\Large \TituloTrabajo}\\[2.5cm]

  \begin{flushright}
    \setlength{\fboxsep}{8pt}
    \fbox{%
      \begin{minipage}{0.4\textwidth}\small
        [PENDIENTE: asignatura].
      \end{minipage}%
    }
  \end{flushright}
  \vspace{1cm}

  {\bfseries \large AUTOR:}\\[0.5em]
  {\large \Autor}\\[2em]
  {\bfseries \large TUTOR:}\\[0.5em]
  {\large \Director}\\[2em]
  {\bfseries \large \Ciudad}\\[1em]
  {\bfseries \large \Fecha}
\end{titlepage}

% ------------------------------------------
%   DECLARACIÓN DE USO DE IA GENERATIVA
% ------------------------------------------
\newpage
\chapter*{Declaración de Delegación a IA Generativa}
\addcontentsline{toc}{chapter}{Declaración de uso de IA Generativa}
Se utilizó ChatGPT/Codex para inspeccionar fuentes, arquitectura y API, consultar referencias y elaborar documentos. No se realizaron solicitudes al backend, pruebas de negocio ni validación docente. La clínica y participación individual están pendientes de registro real.

\noindent\textbf{Declaración enviada por: [PENDIENTE: integrantes reales].}


\chapter{Problema, procesos y arquitectura preliminar}\label{informe-parte1}

\section{1. Revisión de requisitos y flujo
principal}\label{informe-1-revisiuxf3n-de-requisitos-y-flujo-principal}

El problema es organizar la autoevaluación de carreras con referencias
identificables, acceso por ámbito y trazabilidad, evitando cambios sin
historia y preparación con carreras o modelos inadecuados. Fuente:
\path|docs/requirements/spec.md|,
\path|docs/requirements/process-flows.md| y el
\href{../AA2-semana-2/DOSSIER_REQUISITOS_AA2.md}{dossier AA2}.

Se conserva como \textbf{propuesta de incremento de Unidad 1}, pendiente
de ratificación, el acceso, catálogo académico y preparación/consulta de
procesos. La clínica selecciona dos flujos: F3-01, registrar o
corregir una carrera; F3-02, preparar y consultar un proceso.
Inicio/cierre de sesión son soporte transversal. Quedan fuera del
incremento la autorización institucional, activación, carga de
evidencias, valoración, informes, mejora y archivo; que existan
implementados en el producto no los incorpora automáticamente al
ejercicio.

\textbf{Flujo / necesidad }: F3-01 / N-02 y N-05

\textbf{Actor documentado en AA2 }: A-02, administrador funcional con
capacidad de catálogo; lector A-04

\textbf{Requisito e historia }: RF-05--09; HU-02--05

\textbf{Criterios AA2 }: AC-07--17 y 27

\textbf{Operaciones de negocio }: Registrar unidad; registrar carrera;
consultar catálogo; actualizar datos con revisión; registrar decisión de
estado; consultar historial.

\medskip\hrule\medskip

\textbf{Flujo / necesidad }: F3-02 / N-03 y N-04

\textbf{Actor documentado en AA2 }: A-03, gestor autorizado; lector
con permiso aplicable

\textbf{Requisito e historia }: RF-10--12; HU-06

\textbf{Criterios AA2 }: AC-18--24

\textbf{Operaciones de negocio }: Consultar modelo y versión; validar
carrera/período; preparar proceso; consultar estado y referencias.

\medskip\hrule\medskip

\textbf{Flujo / necesidad }: Acceso, integridad y explicación

\textbf{Actor documentado en AA2 }: A-01--04 según operación

\textbf{Requisito e historia }: RF-01--04 y 13--14; RNF-01--03,
05--07 y 10

\textbf{Criterios AA2 }: AC-01--06 y 25--26; VC-01--03, 05--07 y
10

\textbf{Operaciones de negocio }: Autenticar; autorizar; registrar
auditoría; detectar revisión obsoleta; comunicar error seguro.

\medskip\hrule\medskip

\section{2. Análisis textual previo al
modelado}\label{informe-2-anuxe1lisis-textual-previo-al-modelado}

\subsection{F3-01 --- Mantener una carrera
identificable}\label{informe-F3-01--mantener-una-carrera-identificable}

\textbf{Participantes y responsabilidades:} administrador funcional,
proporcionar datos y motivo; frontend, capturar y mostrar resultados;
backend, comprobar identidad, capacidad, reglas y revisión; repositorio,
conservar datos/historia; auditoría, registrar cambio crítico. Frontend,
repositorio y auditoría son participantes técnicos del proceso, no
actores humanos adicionales.

\textbf{Inicio:} solicitar el registro o corrección de una carrera.
\textbf{Precondiciones:} disponer de sesión vigente y
\path|CAREER_ADMIN|; para alta/edición disponer de unidad activa; para
edición obtener detalle y ETag actuales.

\begin{enumerate}
\def\labelenumi{\arabic{enumi}.}
\tightlist
\item
  Capturar código, nombre, unidad y vigencia.
\item
  Validar campos y formato de solicitud.
\item
  Comprobar identidad y capacidad de administración.
\item
  Consultar unidad y verificar su estado activo.
\item
  Verificar unicidad del código normalizado.
\item
  Verificar revisión vigente cuando se solicite modificación.
\item
  Registrar carrera o nueva revisión dentro de la transacción.
\item
  Conservar historial y registrar auditoría del cambio.
\item
  Devolver representación, ETag y, al crear, Location.
\item
  Mostrar resultado y permitir consultar historial.
\end{enumerate}

\textbf{Decisiones:} ¿es alta o modificación?, ¿la entrada es válida?,
¿la unidad está activa?, ¿el código ya existe?, ¿coincide la revisión?
Las actividades de validación de esta descripción no fijan un orden de
precedencia universal de errores: el DTO, filtro y servicio pueden
rechazar antes de etapas posteriores.

\textbf{Alternativa de estado:} consultar revisión; capturar estado y
motivo; validar transición; registrar decisión; conservar carrera e
historia. Desactivar no elimina el recurso. Una unidad con carreras
activas no puede desactivarse. La aplicación devuelve la carrera/unidad
actualizada al registrar la decisión, no una representación separada de
decisión.

\textbf{Excepciones:} 400 por campos/formato; 401 sin sesión; 403 sin
capacidad o CSRF inválido; 404 al consultar identificador inexistente/no
visible; 409 por código duplicado o estado sin cambio; 422 por unidad no
disponible; 428 sin If-Match; 412 por revisión obsoleta. Un fallo
técnico puede resultar en 500; la atomicidad requiere prueba, no se
demuestra por la anotación transaccional sola.

\textbf{Datos:} UUID, código, nombre, unidad, fechas, estado, revisión,
creador y fecha del servidor, motivo y correlación. \textbf{Resultado:}
carrera identificable o rechazo sin declarar cambio completado;
historial y auditoría esperados según AC-14/25. Fuente ejecutable:
\path|CareerController|, \path|CareerAdministrationService|,
\path|Career|, \path|JdbcCareerCatalogRepository| y migración V008, en
las rutas del registro de fuentes.

\subsection{F3-02 --- Preparar y consultar un
proceso}\label{informe-F3-02--preparar-y-consultar-un-proceso}

\textbf{Participantes:} gestor autorizado, seleccionar
carrera/período/modelo; frontend, presentar selecciones; backend,
autorizar y validar; catálogo de carreras, determinar aptitud; catálogo
de modelos, resolver versión publicada; repositorio y auditoría,
conservar proceso y evento.

\textbf{Inicio:} solicitar la preparación de un proceso para una
carrera. \textbf{Precondiciones:} sesión humana vigente,
\path|PROCESS_MANAGE| en el ámbito de carrera, carrera clasificada y
activa en la fecha actual UTC, unidad activa, período contenido en la
vigencia y versión publicada identificada explícitamente.

\begin{enumerate}
\def\labelenumi{\arabic{enumi}.}
\tightlist
\item
  Consultar carreras visibles y seleccionar carrera.
\item
  Consultar modelos/versiones visibles y seleccionar versión exacta.
\item
  Capturar código, nombre y fechas del proceso.
\item
  Comprobar identidad, capacidad y ámbito de carrera.
\item
  Validar actividad de carrera/unidad y pertenencia del período a su
  vigencia.
\item
  Resolver versión publicada sin reemplazarla por otra implícitamente.
\item
  Crear proceso en \path|PREPARATION|, con revisión inicial 0.
\item
  Guardar proceso y registrar auditoría en la transacción.
\item
  Devolver 201, Location, ETag y representación.
\item
  Consultar proceso por identificador y mostrar estado/referencias
  permitidos.
\end{enumerate}

\textbf{Decisiones:} ¿está permitido gestionar esa carrera?, ¿es
utilizable?, ¿pertenece el período a su vigencia?, ¿está publicada la
versión?, ¿es visible el proceso consultado? \textbf{Excepciones:}
entrada ilegible/inválida o versión no publicada \ensuremath{\rightarrow} 400; sin identidad \ensuremath{\rightarrow}
401; gestión fuera de ámbito \ensuremath{\rightarrow} 403; carrera inexistente o proceso no
visible \ensuremath{\rightarrow} 404; carrera inactiva \ensuremath{\rightarrow} 409; período fuera de vigencia \ensuremath{\rightarrow} 422.
No se promete 409 por código de proceso duplicado: debe comprobarse la
regla y su traducción antes de incorporarla al contrato.

\textbf{Datos:} \path|careerId|, \path|code|, \path|name|,
\path|periodStart|, \path|periodEnd|, \path|modelVersionId|; servidor:
ID, número de versión, estado, revisión, actor y fecha.
\textbf{Resultado:} preparación persistida, sin autorización ni
activación automática. \path|ModelApplicabilityService| se consulta en
activación, fuera del recorte: crear en preparación no demuestra
aplicabilidad normativa institucional. Fuente:
\path|SelfAssessmentService.create/get/list| y
\path|SelfAssessmentController|.

\section{3. Arquitectura C4
preliminar}\label{informe-3-arquitectura-c4-preliminar}

El \href{ARQUITECTURA_C4.md}{documento C4} incluye C1, C2, C3,
responsabilidades, tecnologías y fronteras. Los diagramas son vistas
académicas derivadas del código actual, no decisiones nuevas aprobadas.
Se distingue el producto completo del recorte de esta clínica.

En C1 se representan personas y el sistema; no se inventa una
integración activa con CACES, SSO, correo o sistemas académicos. Los
documentos de CACES/UNL son fuentes del dominio, no servicios invocados
por estos recorridos. C2 separa aplicación Angular, backend Spring Boot
y PostgreSQL. El almacenamiento local pertenece al producto completo y
queda fuera de los flujos seleccionados. C3 descompone el backend en
identidad, API, aplicación, dominio, repositorios y puertos de
catálogo/auditoría.

El backend es un \textbf{monolito modular}, según ADR-0001. La
comunicación de navegador a API y de backend a PostgreSQL cruza
procesos; las llamadas entre módulos del backend son internas al mismo
despliegue. No se atribuyen microservicios, broker o transacciones
distribuidas al incremento.

\chapter{Decisiones, recursos y trazabilidad}\label{informe-parte2}

\section{4. Decisiones y alternativas para
discusión}\label{informe-4-decisiones-y-alternativas-para-discusiuxf3n}

\textbf{ID }: DA-01

\textbf{Decisión y evidencia existente }: Monolito modular por
capacidades; ADR-0001

\textbf{Alternativa, ventaja y desventaja }: Microservicios: despliegue
independiente, mayor operación y consistencia distribuida. Capas
globales: arranque simple, más acoplamiento entre capacidades.

\textbf{Justificación / estado }: Conservar límites existentes y
transacciones locales para el recorte. Decisión de producto aceptada; no
nueva aprobación de clínica.

\medskip\hrule\medskip

\textbf{ID }: DA-02

\textbf{Decisión y evidencia existente
}: \path|api/application/domain/infrastructure|; paquetes reales

\textbf{Alternativa, ventaja y desventaja }: Reglas en controlador:
menos clases, acopla HTTP y negocio; CRUD genérico: uniforme,
insuficiente para estados/ámbito.

\textbf{Justificación / estado }: API traduce HTTP; aplicación coordina;
dominio impone invariantes; infraestructura persiste. ArchUnit existe;
no se ejecutó aquí.

\medskip\hrule\medskip

\textbf{ID }: DA-03

\textbf{Decisión y evidencia existente }: Java \textbf{21} en
\path|backend/pom.xml|; Spring Boot \textbf{4.1.1}; Angular
\textbf{22.2.3} en \path|frontend/package.json|

\textbf{Alternativa, ventaja y desventaja }: Cambiar versión de Java:
puede alinear planificación, exige verificar herramientas/CI y
compatibilidad. Mantener configuración: reduce cambio, exige corregir
investigación.

\textbf{Justificación / estado }: \path|docs/architecture/research.md|
menciona Java 25. Documentar divergencia y decidir en equipo; no
certificar compatibilidad mediante lectura del POM.

\medskip\hrule\medskip

\textbf{ID }: DA-04

\textbf{Decisión y evidencia existente }: PostgreSQL como persistencia
principal; Flyway, JDBC y partes JPA

\textbf{Alternativa, ventaja y desventaja }: SQLite auxiliar:
preparación sencilla, no sustituye concurrencia/SQL de aceptación.
Documentos como almacén principal: flexibilidad, menor ajuste a
relaciones/revisiones existentes.

\textbf{Justificación / estado }: Referencias y revisiones requieren
integridad; no introducir event sourcing ni cambiar migraciones
aplicadas. ADR-0018.

\medskip\hrule\medskip

\textbf{ID }: DA-05

\textbf{Decisión y evidencia existente }: Sesión opaca revocable y
autorización por ámbito; ADR-0004

\textbf{Alternativa, ventaja y desventaja }: JWT autocontenido:
verificación local, revocación más compleja. SSO: centraliza identidad,
proveedor/contrato no confirmados.

\textbf{Justificación / estado }: Conservar cookie \path|SEC_SESSION| y
CSRF en mutaciones; no inventar Bearer. Roles técnicos no implican
autoridad funcional.

\medskip\hrule\medskip

\textbf{ID }: DA-06

\textbf{Decisión y evidencia existente }: Puertos
\path|ModelCatalogPort|, \path|AuthorizationPort|, \path|AuditPort|;
adaptadores existentes

\textbf{Alternativa, ventaja y desventaja }: Acceder a tablas de otro
módulo: menor intermediación, rompe límites. Servicio externo
obligatorio: independencia, añade disponibilidad/operación.

\textbf{Justificación / estado }: La consulta de modelo es interna, sin
llamada HTTP externa en el flujo. El MVP no requiere integraciones
externas.

\medskip\hrule\medskip

\textbf{ID }: DA-07

\textbf{Decisión y evidencia existente }: ETag/If-Match y revisiones;
\path|EntityTagSupport|

\textbf{Alternativa, ventaja y desventaja }: Último escritor gana:
sencillo, pierde cambios; bloqueo largo: reduce carreras de
actualización, más contención.

\textbf{Justificación / estado }: Conservar detección de edición
obsoleta. Revisiones iniciales distintas son contrato vigente, no
uniformarlas silenciosamente.

\medskip\hrule\medskip

\textbf{ID }: DA-08

\textbf{Decisión y evidencia existente }: DTO directo y Problem Details
actuales frente al sobre pedido

\textbf{Alternativa, ventaja y desventaja }: Migración coordinada:
satisface forma académica, rompe clientes si se aplica sin versionar.
Mantener Problem Details: compatibilidad, exige acordar excepción a
guía.

\textbf{Justificación / estado }: Presentar sobre objetivo académico y
pedir decisión de compatibilidad. No afirmar que ya está implementado.
Ver contrato.

\medskip\hrule\medskip

Estas alternativas son análisis preparado para discusión, no
alternativas debatidas efectivamente en clase. No se crean ADR aceptados
ni se cambia el estado de los existentes.

\section{5. Recursos, operaciones y
contrato}\label{informe-5-recursos-operaciones-y-contrato}

El \href{CONTRATO_API_INICIAL.md}{contrato inicial} identifica 18
operaciones observadas, sus requisitos, entradas, validaciones,
autorizaciones, resultados y la propuesta de respuesta común. Las rutas
se toman de controladores y la base \path|/api/v1| de
\path|application.yml|; un archivo OpenAPI por sí solo no demuestra
implementación.

Recursos principales: sesión vigente, identidad propia, unidad
académica, carrera, revisión histórica, decisión de estado, modelo,
versión de modelo y proceso. Se utiliza
\path|POST /careers/{id}/state-decisions| como recurso de decisión, ya
existente. No se añade \path|/createCareer|. No se incorpora
\path|/activate| al incremento; su existencia en el producto se declara
como exclusión, sin esconderla.

\section{6. Trazabilidad inicial de arquitectura y
API}\label{informe-6-trazabilidad-inicial-de-arquitectura-y-api}

\textbf{Problema / necesidad }: Acceso autorizado N-01

\textbf{RF / historia / criterio AA2 }: RF-01--04; HU-01;
AC-01--06

\textbf{Proceso }: Acceso transversal

\textbf{Componente principal }: SessionController, SessionService,
ScopeAuthorizationService

\textbf{Recurso / operaciones del contrato }: OP-01--03

\textbf{Verificación posterior }: Sesión válida/inválida, revocación,
ámbito y CSRF.

\medskip\hrule\medskip

\textbf{Problema / necesidad }: Catálogo coherente N-02

\textbf{RF / historia / criterio AA2 }: RF-05--07; HU-02/03;
AC-07--13/27

\textbf{Proceso }: F3-01

\textbf{Componente principal }: CareerController,
CareerAdministrationService, Career

\textbf{Recurso / operaciones del contrato }: OP-04--09, 17

\textbf{Verificación posterior }: Alta, normalización, duplicados,
filtros y visibilidad.

\medskip\hrule\medskip

\textbf{Problema / necesidad }: Corregir sin perder historia N-05

\textbf{RF / historia / criterio AA2 }: RF-08/09; HU-04/05;
AC-14--17; RNF-05

\textbf{Proceso }: F3-01

\textbf{Componente principal }: Servicio de catálogo, repositorio y
auditoría

\textbf{Recurso / operaciones del contrato }: OP-10--12

\textbf{Verificación posterior }: Dos clientes/revisiones, no
sobrescritura, motivo y estado.

\medskip\hrule\medskip

\textbf{Problema / necesidad }: Preparar con referencias exactas N-03

\textbf{RF / historia / criterio AA2 }: RF-10/11; HU-06;
AC-18--22

\textbf{Proceso }: F3-02

\textbf{Componente principal }: SelfAssessmentService, ModelCatalogPort

\textbf{Recurso / operaciones del contrato }: OP-13/14/18

\textbf{Verificación posterior }: Modelo publicado, carrera vigente,
PREPARATION y revisión 0.

\medskip\hrule\medskip

\textbf{Problema / necesidad }: Consultar únicamente lo visible N-04

\textbf{RF / historia / criterio AA2 }: RF-12; HU-06; AC-23/24

\textbf{Proceso }: F3-02

\textbf{Componente principal }: SelfAssessmentService y autorización

\textbf{Recurso / operaciones del contrato }: OP-15/16

\textbf{Verificación posterior }: Listado filtrado e identificador ajeno
oculto.

\medskip\hrule\medskip

\textbf{Problema / necesidad }: Explicar y auditar N-05

\textbf{RF / historia / criterio AA2 }: RF-13/14; AC-25/26;
RNF-06/07/10

\textbf{Proceso }: Ambos

\textbf{Componente principal }: Auditoría, handlers y CorrelationFilter

\textbf{Recurso / operaciones del contrato }: Mutaciones y errores
transversales

\textbf{Verificación posterior }: Atomicidad, correlación y contrato
cliente; pendiente de ejecución.

\medskip\hrule\medskip

Pruebas existentes que se pueden abrir:
\path|backend/src/test/java/ec/edu/unl/sec/architecture/ModuleArchitectureTest.java|;
\path|selfassessment/api/CareerControllerTest.java|,
\path|SelfAssessmentControllerTest.java|;
\path|selfassessment/application/CareerAdministrationServiceTest.java|,
\path|SelfAssessmentServiceTest.java|,
\path|SelfAssessmentAuditServiceTest.java|, bajo el mismo prefijo de
pruebas; \path|tests/acceptance/career-catalog.spec.ts|;
\path|docs/testing/traceability.md|. Esta es vinculación temática por
revisión, no certificación de cobertura individual de todos los AC
nuevos ni ejecución reciente.

\chapter{Clínica y preparación de APE3}\label{informe-parte3}

\section{7. Guion de la clínica y
observaciones}\label{informe-7-guion-de-la-cluxednica-y-observaciones}

Presentar, en este orden: problema/recorte; F3-01/02 y criterios;
C1\ensuremath{\rightarrow}C2\ensuremath{\rightarrow}C3; decisiones DA-01--08; recursos/operaciones; ejemplo 201 y
conflicto de revisión; incompatibilidad del sobre; preguntas pendientes.
Llevar abiertos AA2, análisis textual, C4 y contrato. Cada integrante
debe explicar una decisión y su cadena requisito\ensuremath{\rightarrow}proceso\ensuremath{\rightarrow}componente\ensuremath{\rightarrow}API.

El \href{REGISTRO_OBSERVACIONES.md}{registro de observaciones} separa
hallazgos de esta preparación de retroalimentación auténtica de la
sesión. Para cada observación real se registrarán fuente/fecha,
decisión, cambio o justificación, archivos y evidencia versionada.

\section{8. Cumplimiento y entrega para
APE3}\label{informe-8-cumplimiento-y-entrega-para-ape3}

\textbf{Actividad de la guía }: 1. Revisar requisitos y flujo

\textbf{Evidencia preparada }: Sección 1 y trazabilidad de sección 6

\textbf{Estado / siguiente paso }: Documentado; ratificar
recorte/actores.

\medskip\hrule\medskip

\textbf{Actividad de la guía }: 2. Análisis textual

\textbf{Evidencia preparada }: Sección 2

\textbf{Estado / siguiente paso }: Dos procesos, tareas en infinitivo,
decisiones/excepciones/datos/resultados.

\medskip\hrule\medskip

\textbf{Actividad de la guía }: 3. C4

\textbf{Evidencia preparada }: ARQUITECTURA\_C4.md y diagramas Mermaid

\textbf{Estado / siguiente paso }: Tres niveles; revisar con
equipo/docente.

\medskip\hrule\medskip

\textbf{Actividad de la guía }: 4. Decisiones

\textbf{Evidencia preparada }: Sección 4

\textbf{Estado / siguiente paso }: Ocho decisiones contrastadas y
alternativas por discutir.

\medskip\hrule\medskip

\textbf{Actividad de la guía }: 5--6. Recursos y operaciones

\textbf{Evidencia preparada }: CONTRATO\_API\_INICIAL.md

\textbf{Estado / siguiente paso }: 18 operaciones reales; no nuevos
endpoints ficticios.

\medskip\hrule\medskip

\textbf{Actividad de la guía }: 7. Respuesta común

\textbf{Evidencia preparada }: Contrato, secciones 4--5

\textbf{Estado / siguiente paso }: Sobre académico propuesto;
incompatible con DTO/errores vigentes.

\medskip\hrule\medskip

\textbf{Actividad de la guía }: 8--9. Estados, entradas y validaciones

\textbf{Evidencia preparada }: Contrato, secciones 2--3

\textbf{Estado / siguiente paso }: Revisión estática; confirmar mediante
pruebas de contrato.

\medskip\hrule\medskip

\textbf{Actividad de la guía }: 10. Clínica y ajustes

\textbf{Evidencia preparada }: REGISTRO\_OBSERVACIONES.md

\textbf{Estado / siguiente paso }: Registro preparado; sesión y
retroalimentación pendientes.

\medskip\hrule\medskip


\section{9. Fuentes y límites}\label{informe-9-fuentes-y-luxedmites}

El \href{FUENTES.md}{registro de fuentes} identifica rutas exactas,
decisiones existentes y documentación primaria consultada. Durante la
preparación original ACD3 no se ejecutaron endpoints, pruebas de negocio
ni compilación PDF. Los
ejemplos son didácticos derivados del código, no capturas de ejecución.
Se utilizó ChatGPT/Codex para inspección, análisis y elaboración
documental; la revisión y atribución personal siguen pendientes.

\nocite{c4context,c4container,c4component,rfc9110,rfc5789,rfc6585,rfc9457,fetch,proyecto}
\chapter{Arquitectura C4 preliminar ---
ACD3}\label{c4-arquitectura-c4-preliminar--acd3}

Vistas académicas derivadas del repositorio. Alcance funcional
seleccionado: catálogo y preparación/consulta de procesos con acceso y
auditoría. El producto contiene más módulos, omitidos en C3 para
mantener el foco. Las flechas indican dirección de uso/dependencia, no
flujo temporal.

\section{C1 --- Contexto del sistema}\label{c4-c1--contexto-del-sistema}

\begin{Verbatim}[breaklines=true,breakanywhere=true,fontsize=\small]
flowchart LR
  admin["Persona: administrador funcional\nMantener catálogo autorizado"]
  manager["Persona: gestor autorizado\nPreparar procesos de su ámbito"]
  reader["Persona: lector autorizado\nConsultar información visible"]
  system["Sistema: SIGEAC / SEC-UNL\nOrganizar autoevaluación de carreras\ncon referencias y trazabilidad"]
  admin -->|Registrar y corregir catálogo| system
  manager -->|Preparar y consultar procesos| system
  reader -->|Consultar datos autorizados| system
\end{Verbatim}

\textbf{Frontera:} SIGEAC completo; personas fuera del sistema. A-01
es la condición común de usuario habilitado de estos actores, no un
sistema externo. No se ha identificado dependencia directa de un sistema
externo en los dos flujos. CACES y UNL aportan documentos del dominio;
no se dibujan como APIs. SSO e integraciones de catálogo permanecen
fuera del recorte, sin contrato confirmado.

Fuente: AA2 actores A-01--04; \path|docs/requirements/spec.md|;
\path|docs/adr/0004-opaque-sessions-scope-authorization.md|. Las
competencias institucionales siguen sujetas a confirmación.

\section{C2 --- Contenedores}\label{c4-c2--contenedores}

\begin{Verbatim}[breaklines=true,breakanywhere=true,fontsize=\small]
flowchart LR
  people["Personas autorizadas"]
  subgraph system["Sistema SIGEAC / SEC-UNL"]
    ui["Contenedor: aplicación web\nAngular 22.2.3 / TypeScript\nCapturar entradas y presentar resultados"]
    api["Contenedor: backend API\nJava 21 configurado / Spring Boot 4.1.1\nAutorizar, validar y ejecutar casos de uso"]
    db[("Contenedor: PostgreSQL\nPersistir catálogo, procesos, sesiones y auditoría")]
    files[("Almacén: sistema de archivos local\nEvidencias del producto completo\nFUERA de estos flujos")]
    ui -->|HTTP JSON; cookie y CSRF| api
    api -->|JDBC / acceso SQL; transacciones| db
    api -.->|Puerto de almacenamiento; solo otros flujos| files
  end
  people -->|Usar navegador| ui
\end{Verbatim}

La aplicación web se ejecuta en navegador; el backend es un único
desplegable con módulos por capacidades. El almacén de archivos es un
contenedor de datos en sentido C4, no necesariamente un servicio
independiente ni un contenedor Docker. PostgreSQL es el destino
principal; SQLite es auxiliar de pruebas/desarrollo y no acredita
equivalencia operativa.

\textbf{Contenedor }: Aplicación web

\textbf{Responsabilidad y relación }: Formularios y consultas; cliente
HTTP del backend, sin acceso directo a DB

\textbf{Evidencia }: \path|frontend/package.json|,
\path|frontend/angular.json|, \path|frontend/proxy.conf.json|,
\path|frontend/src/|.

\medskip\hrule\medskip

\textbf{Contenedor }: Backend

\textbf{Responsabilidad y relación }: Autenticación, permisos, reglas,
transacciones y respuesta HTTP

\textbf{Evidencia }: \path|backend/pom.xml|,
\path|backend/src/main/java/ec/edu/unl/sec/|.

\medskip\hrule\medskip

\textbf{Contenedor }: PostgreSQL

\textbf{Responsabilidad y relación }: Persistencia relacional y
revisiones; accesible desde backend

\textbf{Evidencia }: \path|backend/src/main/resources/application.yml|,
\path|backend/src/main/resources/db/migration/|.

\medskip\hrule\medskip

\textbf{Contenedor }: Archivos locales, fuera del recorte

\textbf{Responsabilidad y relación }: Persistencia de evidencias a
través de puerto

\textbf{Evidencia }: \path|docs/adr/0002-file-storage-port.md|,
configuración \path|sec.files.root| en application.yml.

\medskip\hrule\medskip

La relación frontend--API usa HTTP en el proxy local observado; HTTPS es
una condición de despliegue a acordar/verificar, no algo demostrado por
el diagrama. \path|SEC_COOKIE_SECURE| es configurable. No se introducen
colas, cachés o proveedores cloud sin evidencia.

\section{C3 --- Componentes del backend
seleccionado}\label{c4-c3--componentes-del-backend-seleccionado}

\begin{Verbatim}[breaklines=true,breakanywhere=true,fontsize=\small]
flowchart TB
  ui["Contenedor externo al foco: aplicación Angular"]
  db[("Contenedor externo al foco: PostgreSQL")]
  subgraph backend["Contenedor backend Spring Boot — vista parcial"]
    security["Identidad y acceso\nSessionAuthenticationFilter / SessionService\nScopeAuthorizationService"]
    api["API\nCareerController / SelfAssessmentController\nEvaluationModelController / handlers"]
    catalog["Aplicación de catálogo\nCareerAdministrationService"]
    process["Aplicación de procesos\nSelfAssessmentService"]
    domain["Dominio\nCareer / AcademicUnit / SelfAssessmentProcess\nInvariantes y revisiones"]
    model["Puerto de catálogo de modelos\nModelCatalogPort / ModelCatalogAdapter"]
    repos["Puertos de repositorio e implementaciones\nCareerCatalogRepository / SelfAssessmentRepository\nJdbcCareerCatalogRepository / JdbcSelfAssessmentRepository"]
    audit["Auditoría\nAuditPort / SelfAssessmentAuditService"]
    api -->|Invocar casos de catálogo| catalog
    api -->|Invocar casos de proceso| process
    api -->|Consultar modelos vía servicio de modelos| model
    catalog -->|Resolver actor y permisos por puertos| security
    process -->|Resolver actor y permisos por puertos| security
    process -->|Comprobar carrera vigente| catalog
    process -->|Resolver versión publicada| model
    catalog -->|Construir y validar| domain
    process -->|Construir y validar| domain
    catalog -->|Persistir / consultar| repos
    process -->|Persistir / consultar| repos
    catalog -->|Registrar cambio| audit
    process -->|Registrar cambio| audit
  end
  ui -->|HTTP JSON| security
  security -->|Continuar filtro HTTP| api
  repos -->|SQL| db
  audit -->|Persistir evento en DB por infraestructura| db
  model -->|Consultar persistencia de modelos por infraestructura| db
\end{Verbatim}

El bloque de modelos agrupa puerto, adaptador y aplicación para
legibilidad: \path|EvaluationModelController| usa
\path|EvaluationModelService|, mientras \path|SelfAssessmentService|
depende de \path|ModelCatalogPort|. No se afirma que el controlador
inyecte directamente el puerto. El filtro no reemplaza la autorización
de ámbito de cada caso de uso. CorrelationFilter y los handlers son
transversales; no se desdoblan en contenedores.

Dependencias reales a señalar en la clínica:
\path|SelfAssessmentService| llama directamente a
\path|CareerAdministrationService| dentro del módulo selfassessment; no
se inventa un microservicio de carreras ni un puerto que no existe. El
dominio de carreras normaliza código y valida fechas; la aplicación
verifica unidad, permisos, unicidad y revisión; el repositorio conserva
revisiones y usa SQL. Algunas validaciones permanecen en servicios/DTO;
la separación no implica dominio completamente independiente de todo
detalle por el simple nombre de paquete.

\chapter{ACD3 --- Recursos y contrato inicial de
API}\label{contrato-acd3--recursos-y-contrato-inicial-de-api}

\section{1. Estado y reglas de
interpretación}\label{contrato-1-estado-y-reglas-de-interpretaciuxf3n}

\textbf{Base implementada:} \path|/api/v1|, configurada en
\path|backend/src/main/resources/application.yml|. \textbf{Alcance:} 18
operaciones seleccionadas de sesiones, catálogo, modelos y procesos. No
es un inventario completo de la API ni un OpenAPI nuevo. La propuesta de
sobre común de la sección 4 \textbf{no está implementada} en estos
endpoints. Rutas, DTO y estados observados se mantienen separados del
objetivo académico.

Fuentes principales: \path|SessionController|, \path|CareerController|,
\path|EvaluationModelController|, \path|SelfAssessmentController| y sus
servicios/handlers, con rutas exactas en \url{FUENTES.md}. El contrato
v1 está en \path|docs/technical/contracts/openapi.yaml|;
\path|specs/005-complete-evaluation-workflows/contracts/openapi.yaml|
contiene además diseño objetivo v2 y excepciones v1. No se combinan
ambos como si todo estuviera desplegado.

\textbf{Seguridad común S:} cookie de sesión \path|SEC_SESSION|;
comprobar actor humano y capacidad/ámbito; para mutaciones autenticadas,
\path|X-CSRF-TOKEN| vigente. POST de sesión está exceptuado del CSRF de
este filtro. A una lectura protegida sin sesión le corresponde 401; una
capacidad denegada puede producir 403 y un identificador fuera del
ámbito puede ocultarse con 404. El resultado exacto depende del servicio
y orden de rechazo. No se usa \path|Authorization: Bearer| en los
ejemplos.

\textbf{Resultados transversales:} 400 para JSON/UUID/DTO inválidos; 401
falta de autenticación; 403 denegación o CSRF inválido en mutaciones;
500 fallo técnico no controlado. Estas posibilidades no prueban que
todos ocurran en toda operación. 429 está asociado al rechazo de
autenticación por intentos; no se generaliza a cada recurso.

\section{2. Tabla de operaciones
observadas}\label{contrato-2-tabla-de-operaciones-observadas}

Las rutas de la tabla incluyen la base real. Los cuerpos se detallan en
sección 3. \path|GET| no recibe cuerpo en estas operaciones. \path|size|
no existe en el listado de procesos.

\textbf{ID / requisito AA2 }: OP-01 / RF-01

\textbf{Método y ruta }: POST \path|/api/v1/sessions|

\textbf{Entrada y permiso }: Login; sin sesión previa

\textbf{Respuesta de éxito observada }: 201; SessionView; Set-Cookie

\textbf{Rechazos relevantes y causa }: 400 formato; 401 credenciales
rechazadas; 429 límite de intentos.

\medskip\hrule\medskip

\textbf{ID / requisito AA2 }: OP-02 / RF-02

\textbf{Método y ruta }: DELETE \path|/api/v1/sessions|

\textbf{Entrada y permiso }: S; sin cuerpo

\textbf{Respuesta de éxito observada }: 204 sin cuerpo; borrar cookie

\textbf{Rechazos relevantes y causa }: 401 falta de sesión; 403 CSRF.

\medskip\hrule\medskip

\textbf{ID / requisito AA2 }: OP-03 / RF-03

\textbf{Método y ruta }: GET \path|/api/v1/me|

\textbf{Entrada y permiso }: S; sin parámetros

\textbf{Respuesta de éxito observada }: 200; identidad y
funciones/ámbitos vigentes propios

\textbf{Rechazos relevantes y causa }: 401 no autenticado.

\medskip\hrule\medskip

\textbf{ID / requisito AA2 }: OP-04 / RF-05

\textbf{Método y ruta }: POST \path|/api/v1/academic-units|

\textbf{Entrada y permiso }: UnitInput; S y CAREER\_ADMIN

\textbf{Respuesta de éxito observada }: 201; AcademicUnit; Location y
ETag

\textbf{Rechazos relevantes y causa }: 400 entrada; 403 capacidad; 409
código duplicado.

\medskip\hrule\medskip

\textbf{ID / requisito AA2 }: OP-05 / RF-07

\textbf{Método y ruta }: GET \path|/api/v1/academic-units|

\textbf{Entrada y permiso }: page, size, q; S y lectura de catálogo
aplicable

\textbf{Respuesta de éxito observada }: 200; página de unidades visibles

\textbf{Rechazos relevantes y causa }: 400 paginación/búsqueda; 403 sin
capacidad de lectura.

\medskip\hrule\medskip

\textbf{ID / requisito AA2 }: OP-06 / RF-07

\textbf{Método y ruta }: GET \path|/api/v1/academic-units/{id}|

\textbf{Entrada y permiso }: UUID; S y visibilidad de unidad

\textbf{Respuesta de éxito observada }: 200; AcademicUnit; ETag

\textbf{Rechazos relevantes y causa }: 400 UUID; 404 inexistente/no
visible.

\medskip\hrule\medskip

\textbf{ID / requisito AA2 }: OP-07 / RF-06

\textbf{Método y ruta }: POST \path|/api/v1/careers|

\textbf{Entrada y permiso }: CareerInput; S y CAREER\_ADMIN

\textbf{Respuesta de éxito observada }: 201; Career; Location y ETag

\textbf{Rechazos relevantes y causa }: 400 entrada/fechas; 403
capacidad; 409 código duplicado; 422 unidad no disponible.

\medskip\hrule\medskip

\textbf{ID / requisito AA2 }: OP-08 / RF-07

\textbf{Método y ruta }: GET \path|/api/v1/careers|

\textbf{Entrada y permiso }: page, size, q, status, academicUnitId; S y
lectura visible

\textbf{Respuesta de éxito observada }: 200; página de carreras con
total filtrado

\textbf{Rechazos relevantes y causa }: 400 filtros/formato/página; 403
sin lectura.

\medskip\hrule\medskip

\textbf{ID / requisito AA2 }: OP-09 / RF-07/08

\textbf{Método y ruta }: GET \path|/api/v1/careers/{id}|

\textbf{Entrada y permiso }: UUID; S y acceso aplicable

\textbf{Respuesta de éxito observada }: 200; Career; ETag

\textbf{Rechazos relevantes y causa }: 400 UUID; 404 inexistente/no
visible.

\medskip\hrule\medskip

\textbf{ID / requisito AA2 }: OP-10 / RF-08

\textbf{Método y ruta }: PATCH \path|/api/v1/careers/{id}|

\textbf{Entrada y permiso }: UUID, If-Match, CareerInput completo; S y
CAREER\_ADMIN

\textbf{Respuesta de éxito observada }: 200; Career con nueva revisión y
ETag

\textbf{Rechazos relevantes y causa }: 400 entrada/ETag; 404 ID; 409
duplicado; 412 revisión vieja; 422 unidad; 428 falta If-Match.

\medskip\hrule\medskip

\textbf{ID / requisito AA2 }: OP-11 / RF-07/08

\textbf{Método y ruta }: GET \path|/api/v1/careers/{id}/history|

\textbf{Entrada y permiso }: UUID, page, size; S y visibilidad

\textbf{Respuesta de éxito observada }: 200; página de revisiones

\textbf{Rechazos relevantes y causa }: 400 página; 404 carrera no
visible/inexistente.

\medskip\hrule\medskip

\textbf{ID / requisito AA2 }: OP-12 / RF-09

\textbf{Método y ruta }: POST
\path|/api/v1/careers/{id}/state-decisions|

\textbf{Entrada y permiso }: UUID, If-Match, StateChange; S y
CAREER\_ADMIN

\textbf{Respuesta de éxito observada }: 201; Career y ETag; sin Location
explícita en controlador

\textbf{Rechazos relevantes y causa }: 400 DTO/ETag; 404 ID; 409 sin
cambio o reactivación no posible; 412 revisión; 422 motivo/unidad según
escenario; 428 sin precondición.

\medskip\hrule\medskip

\textbf{ID / requisito AA2 }: OP-13 / RF-10

\textbf{Método y ruta }: GET \path|/api/v1/models|

\textbf{Entrada y permiso }: S y MODEL\_READ

\textbf{Respuesta de éxito observada }: 200; lista de modelos

\textbf{Rechazos relevantes y causa }: 401 identidad; 403 capacidad
según servicio.

\medskip\hrule\medskip

\textbf{ID / requisito AA2 }: OP-14 / RF-10

\textbf{Método y ruta }: GET \path|/api/v1/model-versions/{id}|

\textbf{Entrada y permiso }: UUID; S y lectura de modelo

\textbf{Respuesta de éxito observada }: 200; LegacyModelVersion y ETag
cuando la estructura es compatible

\textbf{Rechazos relevantes y causa }: 400 UUID; 404 versión
inexistente/no visible; 409 MODEL\_VERSION\_REQUIRES\_V2 cuando requiere
representación configurada. Consultar no garantiza que esté publicada.

\medskip\hrule\medskip

\textbf{ID / requisito AA2 }: OP-15 / RF-12

\textbf{Método y ruta }: GET \path|/api/v1/processes|

\textbf{Entrada y permiso }: S; filtrar por capacidades/ámbitos; sin
paginación declarada

\textbf{Respuesta de éxito observada }: 200; arreglo de procesos
visibles

\textbf{Rechazos relevantes y causa }: 401 identidad; no prometer
total/paginación.

\medskip\hrule\medskip

\textbf{ID / requisito AA2 }: OP-16 / RF-12

\textbf{Método y ruta }: GET \path|/api/v1/processes/{id}|

\textbf{Entrada y permiso }: UUID; S y acceso a proceso/indicadores
según servicio

\textbf{Respuesta de éxito observada }: 200; SelfAssessmentProcess y
ETag

\textbf{Rechazos relevantes y causa }: 400 UUID; 404 inexistente/no
visible.

\medskip\hrule\medskip

\textbf{ID / requisito AA2 }: OP-17 / RF-04/07

\textbf{Método y ruta }: GET \path|/api/v1/careers/access|

\textbf{Entrada y permiso }: S y acceso al catálogo

\textbf{Respuesta de éxito observada }: 200;
\path|{ "canManage": true }| o false según permiso

\textbf{Rechazos relevantes y causa }: 403 cuenta técnica o sin
capacidad de catálogo. No reemplaza autorización de mutaciones.

\medskip\hrule\medskip

\textbf{ID / requisito AA2 }: OP-18 / RF-11

\textbf{Método y ruta }: POST \path|/api/v1/processes|

\textbf{Entrada y permiso }: ProcessCommand; S y PROCESS\_MANAGE en
carrera

\textbf{Respuesta de éxito observada }: 201; SelfAssessmentProcess;
Location, ETag

\textbf{Rechazos relevantes y causa }: 400 DTO/período invertido/modelo
no publicado; 403 ámbito; 404 carrera ausente; 409 carrera inactiva; 422
período fuera de vigencia.

\medskip\hrule\medskip

RF-13/14 y RNF de privacidad, auditoría, correlación y contrato son
transversales. La existencia de un código en el servicio no demuestra su
respuesta integrada: validar con pruebas/solicitudes en APE3. No se
añade 422 a todo error de validación: los DTO actuales usan 400. No se
asume 409 por cualquier error SQL.

\section{3. Entradas, restricciones y representaciones
actuales}\label{contrato-3-entradas-restricciones-y-representaciones-actuales}

\subsection{3.1 Parámetros y
paginación}\label{contrato-31-paruxe1metros-y-paginaciuxf3n}

Los \path|id| y \path|academicUnitId| son UUID. Catálogos/historial usan
\path|page| desde 0, \path|size| por defecto 20, permitido 1--100;
\path|page| debe ser no negativo y no superar
\path|Integer.MAX_VALUE/100|. \path|q| por defecto vacío, longitud
máxima 300, recortado por el servicio. Carreras admite \path|status|
ACTIVE o DISABLED y filtro de unidad. No inventar filtros de fecha o de
modelo para \path|/processes|.

\path|If-Match| se obtiene de ETag del detalle actual, por ejemplo
\path|"1"|. \path|EntityTagSupport| admite un único entero no negativo
entre comillas fuertes; rechaza forma débil y otras formas con 400,
ausencia con 428. La gramática general de HTTP es más amplia que este
subconjunto del proyecto; revisar esa restricción como decisión de
contrato, no enseñarla como gramática universal. Revisión vieja \ensuremath{\rightarrow} 412
sin sobrescritura.

\subsection{3.2 Cuerpos}\label{contrato-32-cuerpos}

\textbf{Tipo }: LoginRequest

\textbf{Campos y restricciones del DTO / dominio }: username no vacío,
3--100; password no vacío, 12--256

\textbf{Validación de negocio adicional }: Identidad habilitada,
comprobación de credencial, política de intentos; no registrar
contraseña.

\medskip\hrule\medskip

\textbf{Tipo }: UnitInput

\textbf{Campos y restricciones del DTO / dominio }: code no vacío,
máximo 100; name no vacío, máximo 300

\textbf{Validación de negocio adicional }: Código único normalizado;
creador/revisión por servidor.

\medskip\hrule\medskip

\textbf{Tipo }: CareerInput

\textbf{Campos y restricciones del DTO / dominio }: code/name con
límites 100/300; academicUnitId y validFrom obligatorios; validUntil
opcional

\textbf{Validación de negocio adicional }: Unidad activa, código único;
validUntil no anterior a validFrom. Código se recorta, normaliza
espacios y convierte a mayúsculas.

\medskip\hrule\medskip

\textbf{Tipo }: StateChange

\textbf{Campos y restricciones del DTO / dominio }: status obligatorio
ACTIVE/DISABLED; reason no vacío de 10--4000

\textbf{Validación de negocio adicional }: Estado diferente; reactivar
exige unidad activa y clasificación resuelta.

\medskip\hrule\medskip

\textbf{Tipo }: ProcessCommand

\textbf{Campos y restricciones del DTO / dominio
}: careerId/modelVersionId obligatorios; code/name no vacíos, máximo
100/300; periodStart/periodEnd obligatorios

\textbf{Validación de negocio adicional }: Inicio no posterior al fin;
período contenido en vigencia; carrera utilizable en fecha actual UTC;
versión publicada.

\medskip\hrule\medskip

UnitInput, CareerInput y StateChange rechazan campos desconocidos
mediante \path|JsonAnySetter|; no generalizar esta garantía a
ProcessCommand sin comprobar configuración/prueba. Los ID, actor, fecha
y revisión se asignan en servidor.

\subsection{3.3 Formas actuales
relevantes}\label{contrato-33-formas-actuales-relevantes}

\begin{itemize}
\tightlist
\item
  Career: \path|id|, \path|data| con campos de negocio, \path|status|,
  \path|revision|, \path|createdBy|, \path|createdAt|,
  \path|classificationPending|; revisión inicial 1. La entrada de
  creación es plana, no contiene un \path|data| externo.
\item
  AcademicUnit: campos de catálogo propios y revisión inicial 1;
  confirmar schema exacto en su record antes de derivar OpenAPI.
\item
  CatalogPage: \path|items|, \path|page|, \path|size|,
  \path|totalElements|; los totales deben respetar visibilidad.
\item
  SessionView: \path|userId|, \path|expiresAt|, \path|csrfToken|; cookie
  de sesión por Set-Cookie, no contraseña/token opaco en JSON.
\item
  Proceso: \path|id|, \path|careerId|, \path|code|, \path|name|, fechas,
  \path|modelVersionId|, \path|modelVersionNumber|, \path|state|,
  \path|revision|, \path|createdBy|, \path|createdAt|,
  \path|indicators|, \path|authorization|; en preparación inicial,
  revisión 0, indicadores vacíos y autorización ausente.
\item
  Modelo por OP-14: DTO \path|LegacyModelVersion| solo cuando su
  estructura es compatible; en otro caso responde 409
  \path|MODEL_VERSION_REQUIRES_V2|. Para configuración
  completa existe GET \path|/api/v1/model-versions/{id}/configured|. Es
  una ruta real complementaria, fuera de las 18 seleccionadas; no tratar
  DTO legado y configurado como intercambiables. Esta precisión se
  confirmó posteriormente en APE3, casos HTTP-16/16B; no es una
  observación docente.
\item
  Error: Problem Details con \path|type|, \path|title|, \path|status|,
  \path|detail|, \path|instance| y extensiones \path|code|,
  \path|correlationId|, \path|details| según handler. No contiene
  automáticamente \path|success|/\path|errors| del sobre académico.
\end{itemize}

\section{4. Sobre común solicitado: propuesta académica
pendiente}\label{contrato-4-sobre-comuxfan-solicitado-propuesta-acaduxe9mica-pendiente}

La guía solicita esta estructura; se propone como \textbf{contrato
objetivo para discutir}, sin atribuirla al código. Mantener el estado
HTTP real y los headers semánticos. Una respuesta con error no debe usar
200 por llevar \path|success: false|.

\textbf{Escenario }: Alta con representación

\textbf{Uso propuesto del sobre }: success=true; message informativo;
data=representación; errors=null; meta con correlación

\textbf{Estado propuesto / límite }: 201 y Location cuando identifica el
recurso creado.

\medskip\hrule\medskip

\textbf{Escenario }: Lectura individual

\textbf{Uso propuesto del sobre }: data=representación; errors=null

\textbf{Estado propuesto / límite }: 200; ETag donde se implemente
precondición/revisión.

\medskip\hrule\medskip

\textbf{Escenario }: Colección paginada

\textbf{Uso propuesto del sobre}: \texttt{data} = arreglo de ítems; \texttt{meta} = \{page, size, totalElements, correlationId\}

\textbf{Estado propuesto / límite }: 200; no duplicar simultáneamente
paginación en data y meta.

\medskip\hrule\medskip

\textbf{Escenario }: Colección no paginada

\textbf{Uso propuesto del sobre }: data=arreglo; meta con correlación,
sin total/página inventados

\textbf{Estado propuesto / límite }: 200 para procesos mientras no
cambie la operación.

\medskip\hrule\medskip

\textbf{Escenario }: Validación

\textbf{Uso propuesto del sobre }: success=false; data=null; errors con
código/campo/descripción; meta con correlación

\textbf{Estado propuesto / límite }: 400 para DTO inválido; 422 solo
para reglas que lo justifiquen.

\medskip\hrule\medskip

\textbf{Escenario }: Acceso denegado

\textbf{Uso propuesto del sobre }: success=false; data=null; mensaje
seguro; errors con código

\textbf{Estado propuesto / límite }: 401/403 según
autenticación/capacidad; respetar ocultación 404.

\medskip\hrule\medskip

\textbf{Escenario }: Recurso ausente/no visible

\textbf{Uso propuesto del sobre }: success=false; data=null; mensaje
indistinguible para ocultación

\textbf{Estado propuesto / límite }: 404.

\medskip\hrule\medskip

\textbf{Escenario }: Conflicto o precondición

\textbf{Uso propuesto del sobre }: data=null; errors identifican código
estable y acción recuperable

\textbf{Estado propuesto / límite }: 409 para duplicado/estado; 412
revisión vieja; 428 If-Match ausente.

\medskip\hrule\medskip

\textbf{Escenario }: Sin contenido

\textbf{Uso propuesto del sobre }: Ningún sobre ni JSON, ni objeto vacío

\textbf{Estado propuesto / límite }: 204, por ejemplo logout; tampoco se
añade cuerpo a HEAD/304.

\medskip\hrule\medskip

Los campos \path|message|, \path|errors[].field| y
\path|meta.correlationId| son propuesta de forma; no códigos/contratos
aprobados adicionales. \path|Content-Type: application/json| para este
sobre; no etiquetarlo \path|application/problem+json| porque no es un
Problem Details de primer nivel.

\textbf{Compatibilidad que debe resolver el equipo/docente:}

\begin{enumerate}
\def\labelenumi{\arabic{enumi}.}
\tightlist
\item
  Acordar si la entrega exige migrar el producto o documentar un
  contrato futuro y su diferencia actual.
\item
  Si se adopta el sobre, escoger versión/frontera y migrar
  coordinadamente frontend, OpenAPI, handlers, filtros de autenticación,
  pruebas y clientes. No envolver únicamente controladores exitosos y
  dejar errores fuera por accidente.
\item
  Si se conserva Problem Details, registrar aprobación real de la
  excepción académica; no afirmar que el sobre ya está cumplido.
\item
  Mantener cabeceras, permisos y errores semánticos en cualquier
  alternativa. \path|/api/v2| ya aparece como diseño objetivo de 005: no
  reutilizarlo para este cambio sin revisar sus otros compromisos.
\end{enumerate}

\section{5. Ejemplos explicados}\label{contrato-5-ejemplos-explicados}

Valores e identificadores sintéticos; no son respuestas capturadas ni
solicitudes ejecutadas. Las representaciones parciales se etiquetan como
extractos, sin presentarlas como respuesta completa.

\subsection{5.1 Alta de carrera: formato
actual}\label{contrato-51-alta-de-carrera-formato-actual}

\begin{Verbatim}[breaklines=true,breakanywhere=true,fontsize=\small]
POST /api/v1/careers HTTP/1.1
Host: localhost:8080
Content-Type: application/json
Accept: application/json
Cookie: SEC_SESSION=<sesion-sintetica>
X-CSRF-TOKEN: <csrf-de-la-sesion>

{
  "code": "DEMO-C01",
  "name": "Carrera de demostracion",
  "academicUnitId": "11111111-1111-4111-8111-111111111111",
  "validFrom": "2026-01-01",
  "validUntil": null
}
\end{Verbatim}

Si la sesión tiene CAREER\_ADMIN, la unidad es activa y el código es
nuevo, el controlador construye esta clase de respuesta. Las
fechas/identidades siguientes son ilustrativas:

\begin{Verbatim}[breaklines=true,breakanywhere=true,fontsize=\small]
HTTP/1.1 201 Created
Content-Type: application/json
Location: /api/v1/careers/22222222-2222-4222-8222-222222222222
ETag: "1"

{
  "id": "22222222-2222-4222-8222-222222222222",
  "data": {
    "code": "DEMO-C01",
    "name": "Carrera de demostracion",
    "academicUnitId": "11111111-1111-4111-8111-111111111111",
    "validFrom": "2026-01-01",
    "validUntil": null
  },
  "status": "ACTIVE",
  "revision": 1,
  "createdBy": "33333333-3333-4333-8333-333333333333",
  "createdAt": "2026-10-31T12:00:00Z",
  "classificationPending": false
}
\end{Verbatim}

POST solicita crear un elemento de la colección; 201 comunica creación.
Location identifica su URI y ETag permite una edición condicional
posterior. El \path|data| mostrado es un campo real de Career;
\textbf{no} es el sobre de éxito pedido por la guía. Content-Type
describe el JSON enviado y Accept expresa la representación solicitada.

\subsection{5.2 Revisión obsoleta y
recuperación}\label{contrato-52-revisiuxf3n-obsoleta-y-recuperaciuxf3n}

Tras consultar GET \path|/api/v1/careers/{id}| y recibir ETag
\path|"1"|, otro cliente puede guardar revisión 2. El primer cliente
intenta PATCH con \path|If-Match: "1"|, cookie/CSRF y \textbf{todos los
campos de CareerInput}, no solamente name. El servicio rechaza con 412 y
código \path|STALE_REVISION|; debe recargar y conciliar cambios antes de
reenviar. Un error actual didáctico tiene forma:

\begin{Verbatim}[breaklines=true,breakanywhere=true,fontsize=\small]
{
  "type": "urn:sec-unl:problem:stale-revision",
  "title": "Precondition Failed",
  "status": 412,
  "detail": "El recurso cambió; recargue antes de continuar.",
  "instance": "/api/v1/careers/22222222-2222-4222-8222-222222222222",
  "code": "STALE_REVISION",
  "correlationId": "44444444-4444-4444-8444-444444444444",
  "details": {}
}
\end{Verbatim}

La forma se deriva del handler de catálogo; verificar
serialización/cabeceras con prueba integrada antes de presentarla como
captura. \path|ETag| no sustituye autorización ni es token CSRF.

\subsection{5.3 Sobre académico: colección
propuesta}\label{contrato-53-sobre-acaduxe9mico-colecciuxf3n-propuesta}

Para una consulta paginada vacía de carreras, el \textbf{objetivo
propuesto}, no respuesta actual, sería 200 con:

\begin{Verbatim}[breaklines=true,breakanywhere=true,fontsize=\small]
{
  "success": true,
  "message": "Consulta realizada correctamente.",
  "data": [],
  "errors": null,
  "meta": {
    "page": 0,
    "size": 20,
    "totalElements": 0,
    "correlationId": "44444444-4444-4444-8444-444444444444"
  }
}
\end{Verbatim}

Ejemplo propuesto de código duplicado: HTTP 409, Content-Type
application/json:

\begin{Verbatim}[breaklines=true,breakanywhere=true,fontsize=\small]
{
  "success": false,
  "message": "El código ya está registrado.",
  "data": null,
  "errors": [{"code": "DUPLICATE_CODE", "field": "code", "message": "El código ya está registrado."}],
  "meta": {"correlationId": "44444444-4444-4444-8444-444444444444"}
}
\end{Verbatim}

\subsection{5.4 Cierre de sesión}\label{contrato-54-cierre-de-sesiuxf3n}

DELETE \path|/api/v1/sessions| con cookie y CSRF válidos devuelve
\textbf{204 sin cuerpo}, con Set-Cookie para retirar la cookie. No se
devuelve el sobre ni \path|{}|. Reutilizar una sesión revocada en una
ruta protegida debe rechazarse; AC-03 exige comprobarlo, no se afirma
haberlo ejecutado.

\section{6. Métodos, estados y cabeceras: justificación limitada al
caso}\label{contrato-6-muxe9todos-estados-y-cabeceras-justificaciuxf3n-limitada-al-caso}

Según \href{https://www.rfc-editor.org/rfc/rfc9110.html\#section-9}{RFC
9110, §9}, GET consulta, POST procesa contenido según recurso, PUT
reemplaza y DELETE elimina la asociación del recurso. Seguridad
semántica e idempotencia no equivalen a idéntico código en cada
repetición. Aquí DELETE actúa sobre sesión vigente, no borra carreras.
PUT no se selecciona porque no existe reemplazo PUT de estos recursos en
los controladores revisados.

\href{https://www.rfc-editor.org/rfc/rfc5789.html\#section-2}{RFC 5789,
§2} define PATCH para modificaciones parciales. La implementación de
catálogo conserva ID/estado/historial pero exige el bloque completo de
datos editables: documentar esta restricción y discutir si conviene un
documento de cambios específico. No denominarlo JSON Patch ni JSON Merge
Patch: no implementa esos formatos por llamarse PATCH.

200 representa resultado con contenido; 201 creación; 204 éxito sin
contenido; 400 entrada inválida; 401 autenticación necesaria; 403
denegación; 404 ausencia/ocultación; 409 conflicto actual; 412
precondición falsa; 422 contenido no procesable por regla; 500 fallo no
esperado. 428 se fundamenta en
\href{https://www.rfc-editor.org/rfc/rfc6585.html\#section-3}{RFC 6585
§3}, no en RFC 9110. Solo asociar cada estado a escenarios de la tabla,
sin prometer códigos no implementados.

Content-Type y Accept describen/negocian representación; Location
identifica recurso creado; ETag e If-Match coordinan revisión;
Cookie/Set-Cookie mantienen sesión; X-CSRF-TOKEN protege mutaciones
según filtro. 401 requiere revisar la cabecera de desafío
WWW-Authenticate conforme a
\href{https://www.rfc-editor.org/rfc/rfc9110.html\#section-15.5.2}{RFC
9110 §15.5.2}: AA2 ya registró HTTP-AA-01 y el código revisado no la
fija explícitamente. No inventar un esquema Bearer como corrección.

\section{7. CORS y límites
distribuidos}\label{contrato-7-cors-y-luxedmites-distribuidos}

En desarrollo \path|frontend/proxy.conf.json| reenvía \path|/api| a
\path|http://localhost:8080|: el navegador puede usar el mismo origen
visible del frontend. Si se despliegan frontend y API en orígenes
distintos, acordar política CORS, credenciales y exposición de
ETag/Location para que el cliente pueda leerlas. El uso de cookies y
cabeceras de aplicación exige revisar preflight y configuración de
despliegue; no hay evidencia aquí de una política CORS productiva
validada. CORS no concede permisos funcionales ni sustituye CSRF.
Fundamento:
\href{https://fetch.spec.whatwg.org/\#http-cors-protocol}{Fetch,
protocolo CORS}.

El navegador/API y backend/DB son fronteras de comunicación separadas.
No reintentar POST automáticamente tras timeout sin resolver
resultado/idempotencia: no se ha definido clave de idempotencia en este
contrato. Las transacciones entre catálogo, proceso y auditoría son
locales al backend/DB; no se atribuye commit distribuido. Para APE3
habrá que probar pérdida de respuesta, errores, correlación y ausencia
de efectos parciales.
\chapter{ACD3 --- Observaciones y
ajustes}\label{observaciones-acd3--observaciones-y-ajustes}

\section{1. Hallazgos de preparación
documental}\label{observaciones-1-hallazgos-de-preparaciuxf3n-documental}

Estos hallazgos surgen de la inspección asistida del repositorio.
\textbf{No son observaciones recibidas del docente ni de otros equipos.}
Las acciones son propuestas para decidir durante la clínica.

\textbf{ID }: H-01

\textbf{Observación y evidencia }: DTO directos y Problem Details frente
al sobre solicitado; controladores y handlers

\textbf{Decisión propuesta }: Acordar migración compatible o excepción
académica explícita

\textbf{Ajuste documental preparado / pendiente }: Separar contrato
actual/objetivo en CONTRATO\_API\_INICIAL.md. No cambiar backend.

\medskip\hrule\medskip

\textbf{ID }: H-02

\textbf{Observación y evidencia }: Java 25 en research.md frente a
java.version=21 en pom.xml

\textbf{Decisión propuesta }: Confirmar versión de trabajo y corregir la
fuente discrepante tras acuerdo

\textbf{Ajuste documental preparado / pendiente }: C2 muestra Java 21
configurado; DA-03 registra divergencia.

\medskip\hrule\medskip

\textbf{ID }: H-03

\textbf{Observación y evidencia }: PATCH de catálogo exige CareerInput
completo

\textbf{Decisión propuesta }: Acordar semántica/documento de
modificación o mantener restricción explícita

\textbf{Ajuste documental preparado / pendiente }: Documentar entrada
real; no presentarla como JSON Patch.

\medskip\hrule\medskip

\textbf{ID }: H-04

\textbf{Observación y evidencia }: POST state-decisions devuelve carrera
con 201 y sin Location explícita

\textbf{Decisión propuesta }: Aclarar si existe recurso de decisión
consultable y su URI; o revisar estado/representación

\textbf{Ajuste documental preparado / pendiente }: No inventar GET de
decisión ni Location inexistente.

\medskip\hrule\medskip

\textbf{ID }: H-05

\textbf{Observación y evidencia }: Proceso inicia revisión 0; catálogo
en 1; ADR-0018 conserva compatibilidad

\textbf{Decisión propuesta }: Mantener distinción salvo migración
aprobada

\textbf{Ajuste documental preparado / pendiente }: Ejemplos y tabla
conservan ambos valores.

\medskip\hrule\medskip

\textbf{ID }: H-06

\textbf{Observación y evidencia }: Modelo publicado permite preparar;
aplicabilidad se exige después en activación

\textbf{Decisión propuesta }: Confirmar límites de Unidad 1 y validación
normativa

\textbf{Ajuste documental preparado / pendiente }: F3-02 evita
equiparar PREPARATION con autorización/aplicabilidad.

\medskip\hrule\medskip

\textbf{ID }: H-07

\textbf{Observación y evidencia }: AA2 HTTP-AA-01: desafío 401 no fijado
explícitamente en escritor revisado

\textbf{Decisión propuesta }: Comprobar intercambio real y decidir
solución compatible con cookie

\textbf{Ajuste documental preparado / pendiente }: No inventar Bearer ni
declarar conformidad HTTP completa.

\medskip\hrule\medskip

\textbf{ID }: H-08

\textbf{Observación y evidencia }: \path|/processes| devuelve lista sin
parámetros de página

\textbf{Decisión propuesta }: Mantener contrato real o aprobar diseño
paginado

\textbf{Ajuste documental preparado / pendiente }: No fabricar
meta.page/total en contrato actual.

\medskip\hrule\medskip

\section{2. Preguntas para resolver en
clínica}\label{observaciones-2-preguntas-para-resolver-en-cluxednica}

\begin{enumerate}
\def\labelenumi{\arabic{enumi}.}
\tightlist
\item
  ¿Se ratifica el recorte de catálogo y preparación, excluyendo
  activación/gobierno?
\item
  ¿Quién puede administrar catálogo y gestionar procesos en cada ámbito?
  Confirmar asignaciones sin equiparar permisos técnicos con
  competencias institucionales.
\item
  ¿Cómo se cumplirá el sobre de la guía sin romper los clientes v1 y el
  objetivo v2 existente?
\item
  ¿Se conserva PATCH con datos editables completos o se especifica un
  documento parcial distinto?
\item
  ¿Qué recurso crea state-decisions y qué URI/estado deben representar
  la operación?
\item
  ¿Cuál es la versión Java acordada para código, documentación y
  herramientas?
\item
  ¿Qué entorno/orígenes se usarán para probar cookies, TLS, CORS y
  cabeceras?
\item
  ¿Qué alcance y casos de prueba demostrarán atomicidad, visibilidad y
  revisión concurrente?
\end{enumerate}

\section{3. Registro de retroalimentación
auténtica}\label{observaciones-3-registro-de-retroalimentaciuxf3n-autuxe9ntica}

Estado inicial: \textbf{sin registro de sesión aportado}. Completar
exclusivamente con observaciones recibidas realmente.

\textbf{ID }: OBS-01

\textbf{Fecha / persona y rol }: Pendiente

\textbf{Observación recibida }: Pendiente de sesión

\textbf{Decisión y justificación }: Pendiente

\textbf{Cambio realizado o razón para mantener }: Pendiente

\textbf{Requisito / archivo / evidencia }: Pendiente de enlace/commit
real

\textbf{Responsable }: Pendiente

\medskip\hrule\medskip

No convertir H-01--09 en OBS sin una observación real identificada. Si
se mantiene el diseño, registrar el argumento y la fuente, no inventar
un cambio para llenar la tabla.

\section{4. Cambios ya preparados en esta
actividad}\label{observaciones-4-cambios-ya-preparados-en-esta-actividad}

Se elaboraron dos procesos textuales, tres vistas C4, ocho decisiones
con alternativas, selección de 18 operaciones, contrato objetivo
diferenciado y trazabilidad de Semana 2 a componentes/API. Son cambios
en documentación académica; no hubo implementación, migración de
contratos, ejecución de negocio, publicación ni commit nuevo. Las
correcciones derivadas de la clínica permanecen pendientes hasta recibir
retroalimentación real.
\chapter{ACD3 --- Evidencia individual para
completar}\label{individual-acd3--evidencia-individual-para-completar}

\textbf{Estudiante:} pendiente. \textbf{Equipo:} pendiente.
\textbf{Fecha de sesión:} pendiente. \textbf{Rol/aporte efectivamente
realizado:} pendiente. Esta nota no acredita intervención oral ni
revisión personal terminada.

\section{Material preparado para la revisión
personal}\label{individual-material-preparado-para-la-revisiuxf3n-personal}

Seleccionar y revisar al menos una cadena completa de las siguientes;
registrar únicamente lo realizado por el estudiante.

\textbf{Tema disponible }: Catálogo y concurrencia

\textbf{Cadena para explicar }: RF-06/08 \ensuremath{\rightarrow} HU-02/04 \ensuremath{\rightarrow} AC-08--15
\ensuremath{\rightarrow} F3-01 \ensuremath{\rightarrow} CareerAdministrationService \ensuremath{\rightarrow} POST/PATCH careers \ensuremath{\rightarrow}
ETag/412

\textbf{Archivos que abrir }: Dossier AA2, informe ACD3, contrato,
CareerController/CareerAdministrationService/EntityTagSupport.

\medskip\hrule\medskip

\textbf{Tema disponible }: Preparación y límites

\textbf{Cadena para explicar }: RF-10/11 \ensuremath{\rightarrow} HU-06 \ensuremath{\rightarrow} AC-18--22 \ensuremath{\rightarrow}
F3-02 \ensuremath{\rightarrow} SelfAssessmentService \ensuremath{\rightarrow} POST processes \ensuremath{\rightarrow} 201/PREPARATION

\textbf{Archivos que abrir }: Dossier AA2, C4, contrato y servicio de
procesos.

\medskip\hrule\medskip

\textbf{Tema disponible }: Contrato y compatibilidad

\textbf{Cadena para explicar }: RF-14 / RNF-10 \ensuremath{\rightarrow} AC-26 /
VC-10 \ensuremath{\rightarrow} handlers \ensuremath{\rightarrow} Problem Details actual frente a sobre propuesto

\textbf{Archivos que abrir }: Contrato, handlers, ADR-0018 y registro
H-01.

\medskip\hrule\medskip

Las rutas exactas se encuentran en FUENTES.md. Antes de atribuir el
aporte, abrir las fuentes, contrastar el comportamiento y explicar la
decisión con palabras propias.

\section{Registro del aporte
real}\label{individual-registro-del-aporte-real}

\begin{itemize}
\tightlist
\item
  Elementos revisados personalmente: \textbf{pendiente}.
\item
  Hallazgo y fuente concreta (archivo/símbolo/sección):
  \textbf{pendiente}.
\item
  Decisión defendida y alternativa considerada: \textbf{pendiente}.
\item
  Cambio realizado personalmente o justificación para conservar
  propuesta: \textbf{pendiente}.
\item
  Observación realmente recibida, persona/rol y fecha:
  \textbf{pendiente}.
\item
  Uso de la observación y vínculo a OBS real: \textbf{pendiente}.
\item
  Evidencia verificable: \textbf{pendiente de commit, PR, historial o
  acta real}.
\end{itemize}

\section{Ejercicio disponible, no presentado como
ejecución}\label{individual-ejercicio-disponible-no-presentado-como-ejecuciuxf3n}

Explicar por qué PATCH \path|/api/v1/careers/{id}| necesita cookie, CSRF
y If-Match; qué envía CareerInput; por qué una revisión vieja produce
412; qué diferencia existe con 409 por código duplicado; cómo se
recupera sin sobrescribir. Para preparación, explicar por qué 201 y
Location no significan activación ni autorización institucional.

\section{Declaración de
asistencia}\label{individual-declaraciuxf3n-de-asistencia}

Los documentos iniciales fueron preparados con ChatGPT/Codex mediante
inspección estática y consulta de fuentes. Cada integrante debe
identificar su revisión, corrección y participación reales. No
reutilizar el commit base como prueba de contribución a ACD3 ni
presentar el texto generado como registro de una clínica ya realizada.
\chapter{ACD3 --- Fuentes y
evidencia}\label{fuentes-acd3--fuentes-y-evidencia}

Las rutas siguientes son relativas a la raíz del repositorio. Se
inspeccionó código/documentación, sin ejecutar el backend ni pruebas. El
commit base se registra en \path|evidencias/manifest.json|; los
documentos académicos no tienen commit nuevo creado por esta actividad.

\section{Insumos de Semana 2}\label{fuentes-insumos-de-semana-2}

\textbf{Archivo
}: \path|documentos-academicos/APE2-semana-2/specification.md|

\textbf{Uso / límite }: Propuesta de incremento y contratos
preliminares; no aprobación docente.

\medskip\hrule\medskip

\textbf{Archivo
}: \path|documentos-academicos/APE2-semana-2/CONTRATOS_PRELIMINARES.md|

\textbf{Uso / límite }: Línea base de sesiones, catálogo y procesos.

\medskip\hrule\medskip

\textbf{Archivo
}: \path|documentos-academicos/AA2-semana-2/DOSSIER_REQUISITOS_AA2.md|

\textbf{Uso / límite }: Actores, necesidades, HU-01--06, RF-01--14
y RNF-01--10, AC/VC.

\medskip\hrule\medskip

\textbf{Archivo
}: \path|documentos-academicos/AA2-semana-2/MATRIZ_TRAZABILIDAD.md|

\textbf{Uso / límite }: Procedencia documental y pendientes de
validación.

\medskip\hrule\medskip

\textbf{Archivo }: \path|docs/requirements/spec.md|

\textbf{Uso / límite }: Problema, alcance, FR/RNF y clarificaciones de
producto.

\medskip\hrule\medskip

\textbf{Archivo }: \path|docs/requirements/process-flows.md|

\textbf{Uso / límite }: Procesos y límites de responsabilidades.

\medskip\hrule\medskip

\section{Arquitectura y
configuración}\label{fuentes-arquitectura-y-configuraciuxf3n}

\textbf{Archivo / sección }: \path|docs/adr/0001-modular-monolith.md|
--- Decisión/alternativas

\textbf{Uso }: Monolito por capacidades, límites entre módulos.

\medskip\hrule\medskip

\textbf{Archivo / sección }: \path|docs/adr/0002-file-storage-port.md|

\textbf{Uso }: Almacenamiento fuera de DB; solo contexto del producto.

\medskip\hrule\medskip

\textbf{Archivo / sección
}: \path|docs/adr/0004-opaque-sessions-scope-authorization.md|

\textbf{Uso }: Sesiones, hash de token y ámbitos.

\medskip\hrule\medskip

\textbf{Archivo / sección
}: \path|docs/adr/0018-catalog-contract-migration-compatibility.md|

\textbf{Uso }: Compatibilidad v1/v2, catálogo y revisiones, alcance de
capacidades.

\medskip\hrule\medskip

\textbf{Archivo / sección }: \path|docs/architecture/research.md| ---
Versiones base / arquitectura / errores

\textbf{Uso }: Justificaciones existentes; Java 25 difiere del Java 21
configurado. No reiterar versiones como recomendación vigente externa.

\medskip\hrule\medskip

\textbf{Archivo / sección }: \path|backend/pom.xml| ---
parent/properties/dependencies

\textbf{Uso }: Java 21, Spring Boot 4.1.1, dependencias de persistencia,
Flyway y ArchUnit. No prueba entorno instalado.

\medskip\hrule\medskip

\textbf{Archivo / sección }: \path|frontend/package.json| ---
dependencies

\textbf{Uso }: Angular 22.2.3 declarado.

\medskip\hrule\medskip

\textbf{Archivo / sección }: \path|frontend/angular.json| ---
serve/proxyConfig

\textbf{Uso }: Uso de proxy de desarrollo.

\medskip\hrule\medskip

\textbf{Archivo / sección }: \path|frontend/proxy.conf.json|

\textbf{Uso }: \path|/api| hacia localhost:8080; no política CORS
productiva demostrada.

\medskip\hrule\medskip

\textbf{Archivo / sección
}: \path|backend/src/main/resources/application.yml|

\textbf{Uso }: Base API, PostgreSQL, sesiones y almacenamiento.

\medskip\hrule\medskip

\textbf{Archivo / sección
}: \path|backend/src/main/resources/db/migration/V008__career_catalog.sql|

\textbf{Uso }: Esquema/revisiones del catálogo; integridad no ejecutada
aquí.

\medskip\hrule\medskip

\textbf{Archivo / sección
}: \path|docs/technical/contracts/openapi.yaml|

\textbf{Uso }: Contrato común v1 con referencias externas.

\medskip\hrule\medskip

\textbf{Archivo / sección
}: \path|specs/005-complete-evaluation-workflows/contracts/openapi.yaml|

\textbf{Uso }: Objetivo y excepciones por ruta; no demuestra backend v2
disponible.

\medskip\hrule\medskip

\section{Código de los flujos}\label{fuentes-cuxf3digo-de-los-flujos}

Para evitar abreviaturas ambiguas, se registran rutas completas de cada
fuente relevante.

\textbf{Ruta
}: \path|backend/src/main/java/ec/edu/unl/sec/identityaccess/api/SessionController.java|

\textbf{Símbolo y evidencia }: create/revokeCurrentSession/me y DTO;
OP-01--03.

\medskip\hrule\medskip

\textbf{Ruta
}: \path|backend/src/main/java/ec/edu/unl/sec/identityaccess/application/SessionService.java|

\textbf{Símbolo y evidencia }: Creación, autenticación y revocación.

\medskip\hrule\medskip

\textbf{Ruta
}: \path|backend/src/main/java/ec/edu/unl/sec/identityaccess/application/ScopeAuthorizationService.java|

\textbf{Símbolo y evidencia }: Capacidades y ámbitos.

\medskip\hrule\medskip

\textbf{Ruta
}: \path|backend/src/main/java/ec/edu/unl/sec/identityaccess/application/AuthorizationPort.java|

\textbf{Símbolo y evidencia }: Frontera de autorización de casos de uso.

\medskip\hrule\medskip

\textbf{Ruta
}: \path|backend/src/main/java/ec/edu/unl/sec/identityaccess/infrastructure/SessionAuthenticationFilter.java|

\textbf{Símbolo y evidencia }: Autenticación por cookie y CSRF propio en
mutaciones.

\medskip\hrule\medskip

\textbf{Ruta
}: \path|backend/src/main/java/ec/edu/unl/sec/identityaccess/infrastructure/SessionCookie.java|

\textbf{Símbolo y evidencia }: Nombres y atributos de cookie/CSRF.

\medskip\hrule\medskip

\textbf{Ruta
}: \path|backend/src/main/java/ec/edu/unl/sec/identityaccess/infrastructure/AccessProblemWriter.java|

\textbf{Símbolo y evidencia }: Errores anteriores al controlador;
revisar cabecera de desafío 401.

\medskip\hrule\medskip

\textbf{Ruta
}: \path|backend/src/main/java/ec/edu/unl/sec/selfassessment/api/CareerController.java|

\textbf{Símbolo y evidencia }: Rutas, DTO completos, Location y ETag;
OP-04--12/17.

\medskip\hrule\medskip

\textbf{Ruta
}: \path|backend/src/main/java/ec/edu/unl/sec/selfassessment/api/CareerCatalogProblemHandler.java|

\textbf{Símbolo y evidencia }: Conversión de CareerCatalogException a
ProblemDetail.

\medskip\hrule\medskip

\textbf{Ruta
}: \path|backend/src/main/java/ec/edu/unl/sec/selfassessment/application/CareerAdministrationService.java|

\textbf{Símbolo y evidencia }: Autorización, normalización de consulta,
unicidad, unidad, historial y estados.

\medskip\hrule\medskip

\textbf{Ruta
}: \path|backend/src/main/java/ec/edu/unl/sec/selfassessment/domain/Career.java|

\textbf{Símbolo y evidencia }: Datos, vigencia, estado y revisión.

\medskip\hrule\medskip

\textbf{Ruta
}: \path|backend/src/main/java/ec/edu/unl/sec/selfassessment/domain/AcademicUnit.java|

\textbf{Símbolo y evidencia }: Representación de unidad.

\medskip\hrule\medskip

\textbf{Ruta
}: \path|backend/src/main/java/ec/edu/unl/sec/selfassessment/application/CareerCatalogRepository.java|

\textbf{Símbolo y evidencia }: Puerto de catálogo.

\medskip\hrule\medskip

\textbf{Ruta
}: \path|backend/src/main/java/ec/edu/unl/sec/selfassessment/infrastructure/JdbcCareerCatalogRepository.java|

\textbf{Símbolo y evidencia }: Persistencia y visibilidad paginada.

\medskip\hrule\medskip

\textbf{Ruta
}: \path|backend/src/main/java/ec/edu/unl/sec/selfassessment/api/SelfAssessmentController.java|

\textbf{Símbolo y evidencia }: List/get/create y ProcessCommand;
OP-15/16/18.

\medskip\hrule\medskip

\textbf{Ruta
}: \path|backend/src/main/java/ec/edu/unl/sec/selfassessment/api/ProcessExceptionHandler.java|

\textbf{Símbolo y evidencia }: Estado HTTP según Kind; INVALID produce
400.

\medskip\hrule\medskip

\textbf{Ruta
}: \path|backend/src/main/java/ec/edu/unl/sec/selfassessment/application/SelfAssessmentService.java|

\textbf{Símbolo y evidencia }: create/get/list, versión publicada,
PREPARATION, revisión 0, permisos.

\medskip\hrule\medskip

\textbf{Ruta
}: \path|backend/src/main/java/ec/edu/unl/sec/selfassessment/domain/SelfAssessmentProcess.java|

\textbf{Símbolo y evidencia }: Representación y validación de
período/estado.

\medskip\hrule\medskip

\textbf{Ruta
}: \path|backend/src/main/java/ec/edu/unl/sec/selfassessment/application/SelfAssessmentRepository.java|

\textbf{Símbolo y evidencia }: Puerto de procesos.

\medskip\hrule\medskip

\textbf{Ruta
}: \path|backend/src/main/java/ec/edu/unl/sec/selfassessment/infrastructure/JdbcSelfAssessmentRepository.java|

\textbf{Símbolo y evidencia }: Adaptador de persistencia de procesos.

\medskip\hrule\medskip

\textbf{Ruta
}: \path|backend/src/main/java/ec/edu/unl/sec/selfassessment/application/SelfAssessmentAuditService.java|

\textbf{Símbolo y evidencia }: Auditoría del caso de proceso.

\medskip\hrule\medskip

\textbf{Ruta
}: \path|backend/src/main/java/ec/edu/unl/sec/evaluationmodel/api/EvaluationModelController.java|

\textbf{Símbolo y evidencia }: DTO legado/configurado; OP-13/14.

\medskip\hrule\medskip

\textbf{Ruta
}: \path|backend/src/main/java/ec/edu/unl/sec/evaluationmodel/application/EvaluationModelService.java|

\textbf{Símbolo y evidencia }: Lectura de modelos y comprobación de
capacidad.

\medskip\hrule\medskip

\textbf{Ruta
}: \path|backend/src/main/java/ec/edu/unl/sec/evaluationmodel/application/ModelCatalogPort.java|

\textbf{Símbolo y evidencia }: Consulta publicada desde procesos.

\medskip\hrule\medskip

\textbf{Ruta
}: \path|backend/src/main/java/ec/edu/unl/sec/evaluationmodel/infrastructure/ModelCatalogAdapter.java|

\textbf{Símbolo y evidencia }: Adaptador interno, no API externa.

\medskip\hrule\medskip

\textbf{Ruta
}: \path|backend/src/main/java/ec/edu/unl/sec/audit/application/AuditPort.java|

\textbf{Símbolo y evidencia }: Frontera de registro de auditoría.

\medskip\hrule\medskip

\textbf{Ruta
}: \path|backend/src/main/java/ec/edu/unl/sec/shared/api/EntityTagSupport.java|

\textbf{Símbolo y evidencia }: ETag fuerte numérico y errores 400/428.

\medskip\hrule\medskip

\textbf{Ruta
}: \path|backend/src/main/java/ec/edu/unl/sec/shared/api/ApiProblemHandler.java|

\textbf{Símbolo y evidencia }: DTO inválido 400, permisos 403, errores
inesperados 500 y formato de error.

\medskip\hrule\medskip

\textbf{Ruta
}: \path|backend/src/main/java/ec/edu/unl/sec/shared/infrastructure/CorrelationFilter.java|

\textbf{Símbolo y evidencia }: Correlación de solicitud para
errores/auditoría.

\medskip\hrule\medskip

Pruebas a consultar:
\path|backend/src/test/java/ec/edu/unl/sec/architecture/ModuleArchitectureTest.java|,
\path|backend/src/test/java/ec/edu/unl/sec/selfassessment/api/CareerControllerTest.java|,
\path|backend/src/test/java/ec/edu/unl/sec/selfassessment/api/SelfAssessmentControllerTest.java|,
\path|backend/src/test/java/ec/edu/unl/sec/selfassessment/application/CareerAdministrationServiceTest.java|,
\path|backend/src/test/java/ec/edu/unl/sec/selfassessment/application/SelfAssessmentServiceTest.java|,
\path|backend/src/test/java/ec/edu/unl/sec/selfassessment/application/SelfAssessmentAuditServiceTest.java|,
\path|tests/acceptance/career-catalog.spec.ts| y
\path|docs/testing/traceability.md|. Son evidencia de mecanismos
existentes; no se ejecutaron en esta elaboración.

\section{Fuentes técnicas primarias
consultadas}\label{fuentes-fuentes-tuxe9cnicas-primarias-consultadas}

\begin{itemize}
\tightlist
\item
  \href{https://c4model.com/diagrams/system-context}{C4: contexto}:
  frontera de sistema/personas.
\item
  \href{https://c4model.com/diagrams/container}{C4: contenedores}:
  aplicaciones/almacenes y relaciones tecnológicas; no equivale a
  Docker.
\item
  \href{https://c4model.com/diagrams/component}{C4: componentes}:
  descomposición interna de un contenedor.
\item
  \href{https://www.rfc-editor.org/rfc/rfc9110.html}{RFC 9110}: §9
  métodos, §8.8.3 ETag, §10.2.2 Location, §13.1.1 If-Match, §15 estados,
  en particular §15.5.2 para desafío 401.
\item
  \href{https://www.rfc-editor.org/rfc/rfc5789.html}{RFC 5789}: §2,
  semántica de PATCH; no identifica JSON Patch por sí misma.
\item
  \href{https://www.rfc-editor.org/rfc/rfc6585.html}{RFC 6585}: §3
  precondición requerida y §4 límite de solicitudes.
\item
  \href{https://www.rfc-editor.org/rfc/rfc9457.html}{RFC 9457}: §3,
  Problem Details y extensiones; sustento del formato actual, distinto
  del sobre de la guía.
\item
  \href{https://fetch.spec.whatwg.org/\#http-cors-protocol}{WHATWG
  Fetch}: protocolo CORS, credenciales y exposición de cabeceras.
\end{itemize}

Estas fuentes justifican interpretación técnica, no necesidades
institucionales ni aprobación del contrato propuesto. El alcance se
justifica en los documentos de requisitos, no por autoridad de una RFC.

\end{document}
